\chapter{Discussion}


\section{Maneuver Comparison with Previous Study}
Following maneuvers exhibited slightly larger longitudinal gaps and time headways compared to \citeauthor{kutsch_analyzing_2025} (\citeyear{kutsch_analyzing_2025}) (mean gap 11.16 m vs. 9.58 m; mean THW 2.67 s vs. 1.96 s), possibly due to differences in detection thresholds or filtering criteria. For overtaking maneuvers, 573 events were detected here versus 319 previously reported. Despite methodological differences—Kutsch et al. prioritized precise boundary detection while the present approach emphasizes scalability and reproducibility—maneuver durations were similar (mean 7.6 s vs. ~8 s; swerve-out 3.7 s vs. 3.8 s; merge-back 3.8 s vs. 4.17 s), indicating that the current method produces comparable results and is suitable for unsupervised analysis pipelines.



\section{General Clustering Pipeline}
The generalized clustering pipeline successfully captured both local and maneuver-level cycling regimes, demonstrating that a relatively simple and modular framework can reveal meaningful behavioral patterns. While Gaussian Mixture Models and DBSCAN were considered, K-means provided comparable or better separation, with the added benefit of interpretable cluster centroids and computational efficiency. This simplicity prioritizes robustness and interpretability over marginal gains in geometric separation, highlighting that standard unsupervised techniques can suffice for exploratory hierarchical analyses. Future work could explore alternative feature selection methods, dimensionality reduction approaches, or clustering algorithms to assess whether further refinements improve resolution or generalizability, but the current pipeline offers a reproducible baseline.

\section{overtaking behavioral states}
The clustering of overtaking maneuvers revealed two primary behavioral regimes, stable and volatile, consistent with the characterization reported by \citeauthor{bachir_identifying_2025} in free-flow conditions, who observed a similar proportion of volatile windows (~4\%). Scenario differences were pronounced: windows at intersections displayed greater volatility than those on non-intersection segments, likely reflecting increased environmental complexity. Notably, rotation fluctuation contributed strongly to the distinction between regimes at intersections. Unlike \citeauthor{bachir_identifying_2025}, who adjusted this metric to remove curvature effects, the present analysis used the raw signal, which may amplify apparent variability at intersections; accounting for curvature could refine regime separation in future work. 


\section{following behavioral states}
In contrast to overtaking, following behavior exhibited multiple, gradually transitioning regimes rather than sharply distinct clusters. This pattern reflects structured modulation of longitudinal control, consistent with a continuous, adaptive response to surrounding environment rather than event-driven bursts. Scenario-specific differences were also evident: intersection windows tended to show higher variability and additional regimes, suggesting that environmental complexity and traffic interactions influence micro-behavioral adjustments. Compared to overtaking, following maneuvers demonstrate higher regime multiplicity and lower geometric separation, indicating a more progressive and fluid adjustment of riding dynamics. These findings align with the conceptual distinction between event-driven maneuvers, such as overtaking with defined phases (swerve, pass, merge), and continuous behaviors like following, which are largely unbounded except by spatial and temporal constraints. Although cluster separation was moderate, the highly stable partitions and meaningful interpretability indicate that these gradual behavioral states are reproducible and meaningful for modeling longitudinal control patterns in realistic traffic contexts.


\section{maneuver-level aggregation}
The aggregation of local behavioral regimes into maneuver-level representations highlights pronounced structural differences between overtaking and following maneuvers. Following maneuvers typically consist of multiple, shorter regime segments, reflecting ongoing adjustments and continuous modulation of riding dynamics during following interactions. In contrast, overtaking maneuvers are often dominated by a single regime that persists throughout most of the maneuver, producing a cohesive, event-like pattern. These differences persist even when comparing maneuvers of similar duration, indicating that the observed structural patterns are not simply a consequence of maneuver length but instead reflect distinct temporal organization.

Scenario exposure further reinforces this distinction: overtaking is largely concentrated at intersections, consistent with localized, context-dependent execution, whereas following can span multiple segments and exhibits more internal variability. Temporal metrics, including run lengths and centers of mass, indicate that regimes can appear at varying points within maneuvers, particularly in following, highlighting the fragmented and dynamically adaptive nature of these behaviors. Together, these findings underscore a conceptual distinction: overtaking maneuvers are event-driven and temporally cohesive, whereas following maneuvers are progressive, continuous, and responsive to scenarios, providing a structural basis for subsequent clustering of maneuver types in Stage 2.

\section{overtaking maneuver-level execution styles}
The five identified overtaking maneuver clusters reveal how traffic context and rider strategy shape execution styles. Intersection-related overtakes are captured in slow and cautious overtakes (Cluster 0) and accelerating overtakes (Cluster 4). Cluster 0 reflects careful adjustments in constrained intersection environments, with moderate speeds and a balanced mix of stable and volatile regimes. Cluster 4 represents accelerating overtakes, likely occurring immediately after traffic lights or stop-go situations, highlighting bursts of effort triggered by external cues.

In mixed or free-flow conditions, sustained high-speed overtakes (Cluster 1) and smooth high-speed overtakes (Cluster 3) dominate. Cluster 3 corresponds to typical smooth free-flow maneuvers, with stable speeds and minimal variability, representing the “prototypical” passing behavior, while Cluster 1 reflects longer, high-speed overtakes with continuous, controlled acceleration.

Finally, high-steering overtakes (Cluster 2) capture short-duration maneuvers with pronounced lateral motion. Given its low frequency and extreme lateral values, this cluster likely represents rare, constrained situations or residual noise from the overtaking detection algorithm.

Taken together, these clusters illustrate how environmental context (intersection vs. free-flow) and rider strategy (acceleration bursts, steady progression, lateral adjustments) interact to produce distinct overtaking styles. Importantly, this hierarchical, multi-stage unsupervised approach achieves good separation of maneuver-level execution styles. The regime aggregation step plays a crucial role, as cluster differentiation is primarily driven by scenario exposure and dominant micro-regimes, effectively bridging local behavioral states to interpretable maneuver-level patterns.

\section{following maneuver execution styles}
Following maneuvers exhibit substantial internal heterogeneity, reflecting gradual, transient local regimes rather than sharply distinct behavioral states. Maneuver-level clustering reveals this structure: short and quick maneuvers tend to occur at intersections, where lateral variability is higher, likely due to swerving and reactive adjustments to environmental complexity. Longer or spaced maneuvers are concentrated in non-intersection segments, showing steadier lateral control. These patterns demonstrate that scenario exposure partly shapes maneuver execution, but the clustering still captures structured variations within each scenario, highlighting consistent behavioral signatures across environmental contexts.

The identified clusters reflect meaningful differences in interaction intensity with a leading rider. “Catch-up” maneuvers correspond to active acceleration and headway reduction, typically occurring over medium durations and possibly in response to temporary slowing or preparatory for overtaking. “Spaced cruising” and “quick-adjustment” clusters capture varying levels of engagement, from relaxed tracking to reactive micro-adjustments. Moderate and short steady clusters show higher lateral offset variability, indicating responsive adjustments at intersections. These distinctions confirm that the framework can differentiate between low- and high-intensity interaction styles within following maneuvers, which could also imply that these varying styles are scenario dependent.

The successful identification of robust, interpretable clusters despite high internal heterogeneity demonstrates the utility of descriptive, non-generative aggregation of local regimes. By summarizing regime proportions, run lengths, centers of mass, and entropy measures, the method captures structured execution styles across maneuvers. These results support the broader methodological claim that meaningful behavioral patterns can be extracted from hierarchical aggregation of micro-level dynamics, even when transitions between states are gradual rather than sharply delineated.


\section{Infrastructure- and Traffic-Conditioned Discussion}
The conditioning analysis in Stage 3 was designed to characterize the contextual exposure of maneuver execution styles identified in Stage 2. Because clusters were discovered independently of environmental variables, any observed differences in traffic or infrastructure metrics reflect descriptive associations rather than causal influences. This framing emphasizes that Stage 3 examines \textit{where} certain execution styles tend to occur, not \textit{why} they emerge.

Within both intersection- and non-intersection–dominant clusters, infrastructure variables such as lane width and lane type showed minimal variation. This confirms that differences in execution styles are not driven by geometric features within the same scenario class, allowing interpretation to focus on traffic exposure metrics.

Overtaking maneuvers exhibited moderate associations with environmental conditions. Accelerating intersection overtakes were generally accompanied by higher pedestrian conflict indicators (DRAC and TTC), whereas slow intersection overtakes occurred under lower exposure. Non-intersection overtakes were largely similar across clusters, with minor differences in bicycle counts. The small sharp-steering overtakes cluster (n = 16) produced stepwise ECDF curves, limiting its interpretability. These observations suggest that overtaking behavior may partially adapt to surrounding traffic demands, though the effects are associative rather than causal.

In contrast, following maneuvers showed limited differentiation across clusters once scenario context was controlled. ECDF distributions largely overlapped, and the ordering across clusters was inconsistent across metrics and participant types. This indicates that following styles are primarily driven by leader–follower interaction dynamics, including gap regulation, speed matching, and response timing, rather than by broader environmental conditions. This finding aligns with Stage 2 observations, where following clusters were more gradual and less distinct than overtaking clusters, highlighting the intrinsic nature of interaction-driven behavior.

The scenario-based separation applied in Stage 3 effectively controlled for environmental context, such as intersection versus non-intersection segments. This separation clarifies where maneuvers tend to occur, suggesting that most environmental variation is captured by the scenario split. In contrast, the intrinsic execution characteristics of each maneuver — including speed profiles, acceleration patterns, lateral movements, and regime dynamics — describe how maneuvers are performed. Within scenario-dominant clusters, these execution styles showed minimal differentiation across traffic exposure metrics, indicating that the internal structure of maneuvers appears largely independent of the immediate environment. In other words, the scenario split accounts for the location and context of maneuvers, while execution style within those contexts seems to reflect intrinsic behavioral strategies, a pattern that is particularly evident for following maneuvers.




In this study, overtaking maneuvers were captured as independent events by stationary drones, without the ability to track individual cyclists across multiple occurrences. As a result, clustering analysis of overtaking cannot reveal persistent cyclist-specific behavior (“styles”), but only transient maneuver types. Our results confirm this: clusters reliably distinguish between active, interaction-driven overtakes and passive, free-flow overtakes, yet these clusters reflect event-level modes rather than enduring behavioral regimes. In contrast, following behavior is continuous within a trajectory and can be clustered into regimes, such as “close, reactive” versus “far, stable” following, because it exhibits persistent patterns over time. This distinction highlights a fundamental methodological limitation: regime-level analyses require repeated observations of the same agent, which is not available for overtaking in the current dataset.



For volatile riding regimes, the ECDF of $distance_traffic_light$ exhibits a pronounced knee at approximately 20 m. 
This transition aligns with prior findings indicating that cyclist behavior is influenced by traffic lights within a ~20 m range, suggesting that this distance marks a behaviorally relevant regime boundary for volatile riders. 
visual attention and environmental scanning intensify as cyclists enter <35m from the signal \cite{rupi_visual_2019}, aswell as cyclists tend to decelerate/accelerate within 30 m of a stop line at signalised intersections \cite{haifeng_research_2013}



resulting maneuver clusters: these are modes, not styles. for a style analysis the focus has to be on single cyclists and how they behave across multiple maneuvers. but this framework might be a first step in this direction.

The weak separability of local regimes in following scenarios suggests that following behavior may be better described by continuous relational dynamics rather than discrete riding regimes.


extend framework by:
 - clustering traffic state as done by \citeauthor{ren_lane-change_2025, hu_high-resolution_2022} 
 - different window segmentation approach: acc/decel \citeauthor{yarlagadda_exploring_2022}, 
 - frequency-based features on driving parameters \cite{chen_feature_2023}
 - exploiting sequenciality of windows (currently ordering is only partially captured)
 - differentiating different maneuver phases (pre-/post-overtake)
 
 
\subsection{Linking Behavioral Clusters to Traffic Exposure}
Stage 3 analyses reveal systematic relationships between maneuver clusters, local regimes, and traffic exposure.

For overtaking maneuvers, slow-steady overtakes (Cluster 0) combine mixed stable and volatile regimes and are associated with low DRAC values and limited participant exposure, suggesting cautious adjustments in constrained intersections. Accelerating overtakes (Cluster 4) maintain stable regimes but encounter higher DRAC and lower pedestrian min TTC, reflecting more frequent proximity events. Fast-wide overtakes (Cluster 1) experience higher bicycle counts and elevated DRAC, consistent with overtaking multiple cyclists or denser traffic, while fast-smooth overtakes (Cluster 3) occur under lower traffic exposure, matching their stable, streamlined behavior. Sharp-steering overtakes (Cluster 2) show mixed regimes and sparse ECDF patterns, indicating agile maneuvers under variable exposure.

For following maneuvers, intersection-dominant clusters (0, 2, 4) show minimal differences in traffic metrics, suggesting consistent behavior across clusters. Non-intersection clusters reveal meaningful contrasts: prolonged close following (Cluster 3) maintains stable regimes over longer durations and experiences more low TTC events with pedestrians, while spaced cruise (Cluster 5) also exhibits stable regimes but with higher min TTC and reduced interactions. Cluster 1 diverges across metrics, reflecting heterogeneous scenarios and flexible adjustment strategies.